\begin{recipebox}{Chicken Pad Thai}{2--3}{30 min}
\begin{ingredients}
  \item \textbf{Sauce:}
  \begin{itemize}[leftmargin=2em]
    \item 4 Tbsp fish sauce
    \item 2 Tbsp brown sugar
    \item 3 Tbsp rice wine vinegar
    \item 3 Tbsp soy sauce
  \end{itemize}
  \item \textbf{Main:}
  \begin{itemize}[leftmargin=2em]
    \item Drizzle of oil
    \item 2 boneless, skinless chicken thighs, thinly sliced
    \item 1 clove garlic, finely chopped
    \item Pinch of chilli flakes
    \item 1 egg, beaten
    \item 100 g mung bean sprouts
    \item 200 g rice noodles
    \item 2 spring onions, thinly sliced
    \item 2 Tbsp crushed peanuts
    \item Lemon wedges, to serve
  \end{itemize}
\end{ingredients}

\begin{method}
  \item In a small bowl, mix together the sauce ingredients. Set aside.
  \item Heat oil in a large pan or wok. Add chicken and cook until golden and cooked through. Add garlic and chilli flakes, stirring for 1 minute.  
  \item Push chicken to the side and pour the beaten egg into the centre. Scramble gently until cooked through, then mix with the chicken.  
  \item Soak rice noodles as per packet instructions (usually 2–3 minutes to soften). Drain well in a colander.  
  \item Add noodles to the pan along with the prepared sauce, bean sprouts, spring onions, and peanuts. Toss everything together and simmer for a few minutes until the sauce reduces and noodles are tender.  
  \item Serve immediately, topped with extra peanuts, spring onions, and lemon wedges.  
\end{method}

\begin{notes}
You can swap chicken for prawns, tofu, or beef.  
Adjust the chilli flakes to taste for more or less heat.  
Traditional Pad Thai often includes tamarind paste in the sauce — add 1 Tbsp if you have it for extra tang.  
\end{notes}
\end{recipebox}